\documentclass[10pt,a4paper]{amsart}
\usepackage[margin=19mm]{geometry}
\usepackage{amsmath,amssymb,mathtools}
\usepackage[scr=rsfs,cal=euler]{mathalfa}
\usepackage{xfrac}
\usepackage[unicode,hidelinks,bookmarks=false]{hyperref}
\newcommand{\ie}{i.e.\ }
\newcommand{\cf}{cf.\ }
\newcommand{\resp}{respectively\ }
\newcommand{\Subsection}{\subsection}
\providecommand{\nobreakdash}{\nobreak\hbox{-}}
\setlength{\emergencystretch}{2em}
\multlinegap0pt
\def\Q{{\mathbb Q}}
\def\F{{\mathbb F}}
\usepackage{amssymb}
\usepackage{xcolor}
\newtheorem{thm}{Theorem}
\title{Modular curves and bad reduction}
\author{Adam Logan, David McKinnon}
\date{Source: arXiv:2604.09536v1 (CC BY 4.0)}
\begin{document}
\maketitle

\noindent\emph{Source and licence.} Modular curves and bad reduction. Version arXiv:2604.09536v1 (CC BY 4.0), distributed by arXiv under the Creative Commons Attribution 4.0 International licence, http://creativecommons.org/licenses/by/4.0/, as recorded in the arXiv licence metadata for this exact version and shown on the abstract page https://arxiv.org/abs/2604.09536v1. Full licence notice: \texttt{licenses/arxiv-2604.09536-CC-BY-4.0.txt}.\par\smallskip

\noindent\textit{Editorial scope.} Counted unit: the article's thm thm:stronger (source0.tex lines 165--174). The article's complete original proof of it is delivered with it (source0.tex lines 176--185, one \texttt{proof} environment of 10 lines). No mathematics is written, repaired or re-derived: the statement and the proof are the article's own lines, in the article's own notation, and no retained line is substituted or re-flowed.  No further source-local context is needed: every cross-reference of every retained line resolves inside the two delivered spans.  Nothing was omitted from any retained span, so the package carries no omission record.  No retained line contains a citation, so the wrapper carries no bibliography and no bibliography entry.  No retained line contains a figure environment, an image-inclusion command or drawing code, so no image is delivered and the package declares no asset dependency.  Editorial formatting only, in addition: an AMS \texttt{amsart} preamble that restates the article's own declarations and macros used by the retained lines, namely the environments \texttt{thm} on separate counters, together with the standard packages those lines need.  The retained environments print in this document as Theorem 1 in order of appearance, whereas the article prints the same items with the numbering of its own full document.  The complete native source of the article as distributed in the arXiv e-print for this exact version is bundled unchanged as \texttt{sources/198217/source0.tex}.
\begin{thm}\label{thm:stronger}
  Let $E$ be a modular elliptic curve as before, let $K$ be a quadratic
  field, let $p$ be a prime inert in $K$, and suppose that every $\F_p$-point
  of $E$ is represented by a cusp and that the map $E \to \P^1$ to the
  $x$-line is surjective on $\F_p$-points.  Let $E^\sigma$ be the twist of
  $E$ corresponding to $K$.  Suppose that the natural map
  $E(\Q) \oplus E^\sigma(\Q) \to E(K)$ is surjective.  Then every elliptic
  curve over $K$ represented by a $K$-point of $E$ does not have potentially
  good reduction at $p$.
\end{thm}

\begin{proof} Let $P \in E(K)$ and write $P = C + R$, where $C \in E(\Q)$
  %is congruent mod $p$ to a cusp  % I don't see why this needs to be mentioned here
  and $R \in E^\sigma(\Q)$.  Let $R_0$ be the
  image of $R$ in the $x$-line: then $R_0$ reduced mod $p$ is also the image of
  a $\Q$-point of $E$, since the reduction map is surjective on $\F_p$-points.
  This is only possible if $R_0$ is a Weierstrass point mod $p$, so 
  $R \bmod p \in E(\F_p)$ and hence it is the reduction of a point of $E(\Q)$.
  Thus $P \bmod p$ is the reduction of a point of $E(\Q)$ as well; by hypothesis
  it is the reduction of a cusp.  The argument proceeds as before.
\end{proof}

\end{document}
